% Ad Familiares 11.16 --- headnote
% ad-familiares-11-16, 15 July 43 BC, Rome

\textit{Cicero to D. Brutus, consul-designate, from Rome on
15 July 43 BC --- the Perseus dateline carries the wide
range \textit{Scr. Romae inter ex. m. Apr. et in. Quint. a.
711 (43)}, but the day after 11.22 (14 July) is the
conventional placing, and matches the manuscripts' running
order. The subject is the praetor's elections, and the
particular candidate is L. Aelius Lamia, a leading
\textit{eques} who in 58 BC --- as Cicero recounts in
section 2 --- was driven out of Rome by the consul Gabinius
for resisting Cicero's exile, the first Roman citizen so
treated at Rome.}

\textit{Form and substance are exquisitely matched. The
letter opens with an unusually delicate parable about the
right moment for handing in a letter --- Cicero has briefed
his courier to watch for it --- and the cover for that
delicacy is the request that follows: that Brutus, as
president of the equestrian centuries in the upcoming
elections, throw their weight (via Lupus) behind Lamia.
The pivot at the end of section 2 is the boldest stroke:
``persuade yourself of this, my Brutus: I am the one
standing for the praetorship.'' This is the last letter
to D. Brutus that has come down to us. Within five months
both correspondents will be dead.}
